\section{Methods}

We formalize the optimization of audio restoration operators within a differentiable framework that unifies signal processing priors with deep learning flexibility. Let $\mathbf{x}$ denote an observed noisy waveform and $\hat{\mathbf{x}}$ its restored counterpart, where both are sampled at $F_s = 48\,\text{kHz}$. The core objective is to minimize the residual error between a target clean reference $\mathbf{y}_{\mathrm{id}}$ (obtained via high-fidelity wavetable synthesis or ground-truth recordings) and the output of a parametric operator $\Theta(\cdot)$, parameterized by weights $w_\theta$, architecture hyperparameters $h$, and latent control variables. We define the primary score as the normalized residual ratio:
\begin{equation}
R = \text{clamp}\left(1 - \frac{\|\hat{\mathbf{x}} - \mathbf{y}_{\mathrm{id}}\|_2 / F_s}{\max(\|\mathbf{y}_{\mathrm{id}}\|_2/F_s, \epsilon)}, 0, 1\right),
\end{equation}
where $\epsilon = 10^{-6}$ prevents division by zero. This metric $R$ aligns with the principle that perfect restoration yields a score of unity, while noise leakage reduces it proportionally to the residual root-mean-square relative to the ideal signal energy \citep{jaderberg2017pbt}. Unlike standard mean-squared error losses which may favor trivial solutions in unsupervised settings, our formulation explicitly penalizes deviations from the known spectral envelope of clean audio.

The operator $\Theta$ is constructed using Differentiable Digital Signal Processing (DDSP) primitives \citep{engel2020ddsp}, allowing gradients to flow through time-domain operations such as wavetable synthesis and overlap-add convolution. This architecture contrasts with purely end-to-end waveform models like WaveGrad, which lack explicit control over harmonic structure or transient preservation \citep{shan2021diffwave}. To address the hyperparameter sensitivity inherent in deep learning systems, we employ Population Based Training (PBT) rather than static grid search or manual tuning. PBT maintains a population of agents where each agent trains on a subset of data and periodically inherits successful configurations from outperforming peers \citep{jaderberg2017pbt}. This asynchronous strategy effectively utilizes the fixed computational budget to jointly optimize model weights, architecture topology, and learning rates without requiring explicit architectural priors.

Our search procedure operates in three distinct branches: Reinforcement Learning (RL), Neural Architecture Search (NAS), and a hybrid Combo branch that interleaves gradient-based updates with evolutionary selection \citep{elsken2019nas}. In the RL branch, we utilize Proximal Policy Optimization to maximize expected $R$ over a validation set of corrupted wavetables. The policy network outputs parameters for DDSP modules (e.g., modulation indices, attack times) conditioned on input features extracted by a shared encoder. For NAS and Combo branches, we relax discrete architectural choices—such as the number of synthesis stages or filter bank resolutions—to continuous variables via DARTS-style operations \citep{liu2019darts}, enabling gradient-based exploration of the architecture space. The Combo branch further incorporates meta-learning techniques inspired by Model-Agnostic Meta-Learning (MAML) to accelerate adaptation across diverse noise types, though we prioritize direct optimization for this specific restoration task where clean references are available \citep{finn2017maml}.

Implementation constraints significantly affect validity. Training occurs on NVIDIA GeForce RTX 3090 GPUs with mixed-precision arithmetic to manage the high-dimensional state spaces typical of audio synthesis models. We enforce a maximum training duration of $168$ hours per branch, sampling history every iteration to track convergence dynamics. Crucially, we avoid data augmentation strategies that introduce artifacts not present in real-world recordings; instead, noise corruption is simulated by adding colored Gaussian noise and impulse distortions derived from physical crackle models \citep{lehtinen2018n2n}. The search space includes latent variables governing wavetable wrap behavior to mitigate the "crackle" phenomenon often observed when periodic waveforms are truncated or aliased. By grounding our approach in both statistical signal reconstruction principles and modern automated design paradigms, we aim to discover operators that generalize across unseen noise profiles while maintaining fidelity to perceptual audio quality metrics.
